\documentclass[10pt]{article}
\usepackage[margin=20mm]{geometry}
\usepackage{amsmath,amssymb,mathtools,fontspec,unicode-math,xurl,hyperref}
\IfFontExistsTF{DejaVu Serif}{\setmainfont{DejaVu Serif}}{\setmainfont{Latin Modern Roman}}
\setmathfont{latinmodern-math.otf}
\setlength{\parindent}{0pt}\setlength{\parskip}{6pt}\sloppy
\title{Scientific Review of
Spherical Caps do Not Always Maximize
Neumann Eigenvalues on the Sphere}
\author{AI-assisted technical review}
\date{}
\begin{document}
\maketitle
Abstract: The manuscript proves that sufficiently large spherical caps fail to maximize the first nonconstant Neumann eigenvalue at fixed area, using a symmetric four-hole perturbation. This review identifies one confirmed scope defect: the general hole hypotheses in Section 3 do not exclude a bounded complementary component. A translated annulus invalidates Lemma 7, an auxiliary exterior problem, and the fixed-index conclusion of Proposition 3. The actual main construction uses filled ellipses, and Section 6 expressly assumes a simply connected hole. The demonstrated defect is classified as Minor; no failure of Theorem 1 is established. The numerical helmet-domain discussion is assessed within its expressly non-rigorous scope.\par
Paper and version. M25-0392: Dorin Bucur, Richard S. Laugesen, Eloi Martinet, and Mickaël Nahon, Spherical Caps do Not Always Maximize Neumann Eigenvalues on the Sphere. Exact object reviewed: arXiv:2503.15385v1, 19 March 2025, 23 pages. Frozen journal identity: Geometric and Functional Analysis (2025), DOI `10.1007/s00039-025-00721-1`. The conclusions apply to this author manuscript; equivalence to the journal version has not been established. Physical and printed page numbers agree.\par
\section*{Scope and principal argument}
All 23 pages of the selected manuscript were read. The review examined the variational indexing, eigenvalue-cluster hypotheses, boundary normal conventions, ansatz and residual estimates in Sections 3 and 7, parameter choices proving Theorem 1, numerical status of Section 5, and fixed-shape restriction in Section 6. Critical formula pages and Figures 1–4 were also inspected as actual rendered source images. Source and evidence byte counts and SHA256 bindings accompany the report.\par
Theorem 1 asserts the existence of \(\delta>0\) such that for every \(\theta\in(\Theta-\delta,\pi)\) there is an open set of the same area as \(B_\theta\) with strictly larger \(\mu_1\). Here \(\mu_1(B_\Theta)\sin^2\Theta=1\). Proposition 2 supplies the double eigenspace generated by \(g(t)\cos\varphi\) and \(g(t)\sin\varphi\) and the hot-spot structure of \(g\). These are cited results from Langford and Laugesen [LL]. Their publisher-visible appendices were consulted, including the Wronskian comparison and Legendre-function argument. The complete proofs of their principal Propositions 3.1, 4.1 and 4.2 were not accessible in the material obtained. This review checks the deduction from the manuscript’s stated Proposition 2 and does not claim a complete independent reproof of that external input.\par
Equation reference: ellipse. For the filled ellipses in Section 4, the virtual-mass coefficients are \[
 \frac{(M_{\omega_\lambda})_{11}}{\pi}=\frac{2}{1+\lambda^{-2}},
 \qquad
 \frac{(M_{\omega_\lambda})_{22}}{\pi}=\frac{2}{1-\lambda^{-2}}.\] The fourfold rotational symmetry preserves the double cluster. On the limiting two-dimensional real eigenspace a quarter-turn has no real one-dimensional invariant subspace, and a commuting self-adjoint operator is scalar there. Continuity of the isolated spectral representation gives the same conclusion for sufficiently small holes.\par
Equation reference: endpoint. For \(\theta>\Theta\), placing the holes at \(t_{\max}\) eliminates the \(g'(t)^2\) term. The remaining terms become positive for \(\lambda>1\) sufficiently close to one. At \(\theta=\Theta\), the expansion on page 11, with \(t=\Theta-\tau\) and \(\lambda=1+\eta\), is \[
 (-2\tau\cot\Theta-\eta)\frac{g(\Theta)^2}{\sin^2\Theta}
 +O\!\left(\tau^2+\eta^2+\frac{\tau^4}{\eta}\right).\] Choosing \(\eta=-\tau\cot\Theta>0\) makes the leading term positive, of order \(\eta\), with remainder \(O(\eta^2)\). The parameters of each fixed ellipse are chosen before \(\varepsilon\) tends to zero. No uniform remainder over arbitrarily degenerating ellipses is needed for this existence argument. Continuity gives a left neighborhood of \(\Theta\). The issue below does not supply a counterexample to these filled-ellipse steps.\par
\section*{A missing topological hypothesis}
Problem 1 (Minor: bounded components of the hole complement).  Section 3.1, pages 3–4, permits a smoothly bounded closed set \(\omega\subset\mathbb R^2\) whose interior contains the origin, without requiring \(\mathbb R^2\setminus\omega\) to be connected. Under these stated assumptions, the harmonic-extension estimate in Lemma 7 (page 22) and the fixed-index expansion in Proposition 3 (page 5) are false. The asserted matrix-valued exterior corrector on page 4 may not exist.\par
\subsection*{Exact admissible geometry}
Equation reference: annulus. Let \(c=(2,0)\) and take \[
 \omega=\{x\in\mathbb R^2:1\leq |x-c|\leq3\}.\] This closed annulus is compact and connected, with two smooth boundary circles. Since \(1<|c|=2<3\), its interior strictly contains the origin. The removed set itself is connected; its complement has the bounded component \(D_1(c)\) and the unbounded component \(\mathbb R^2\setminus\overline{D_3(c)}\).\par
\subsection*{Failure of Lemma 7}
Equation reference: h, energy. For the open set in Lemma 7 use the interior of [annulus] (\#annulus), write \(r=|x-c|\), and define \[
 h(x)=
 \begin{cases}
 1,&r\leq1,\\
 \log(3/r)/\log3,&1<r<3,\\
 0,&r\geq3.
 \end{cases}\] The traces match at both circles, so \(h\) is a compactly supported \(H^1(\mathbb R^2)\) function. It is harmonic in the annulus, and \[
 \int_{\omega}|\nabla h|^2\,\mathrm dx
 =\frac{2\pi}{(\log3)^2}\int_1^3\frac{\,\mathrm dr}{r}
 =\frac{2\pi}{\log3}>0,\qquad
 \int_{\mathbb R^2\setminus\omega}|\nabla h|^2\,\mathrm dx=0.\] No finite constant in Lemma 7 satisfies the asserted estimate. The single-mean Poincaré inequality after inversion is the failed step: functions constant at different values on distinct components have zero gradient and need not equal one common mean.\par
\subsection*{Failure of the fixed-index spectral assertion}
Equation reference: chart, zero. The same geometry satisfies every stated geometric hypothesis preceding Proposition 3. Set \(M=\Omega=\mathbb S^2\), use one marked point at the south pole, and let the boundary vector field be zero. Empty boundary is explicitly allowed. A conformal chart with conformal factor one at the origin is \[
 \psi(x)=\left(\frac{4x_1}{4+|x|^2},
              \frac{4x_2}{4+|x|^2},
              \frac{|x|^2-4}{4+|x|^2}\right).\] For a fixed chart ball \(B_r\) and \(0<\varepsilon<r/5\), the entire set \(\varepsilon\omega\) lies inside \(B_r\). The cluster \(I=\{1,2,3\}\) satisfies the required gap conditions: \[\mu_0(\mathbb S^2)=0<\mu_1(\mathbb S^2)=\mu_2(\mathbb S^2)
 =\mu_3(\mathbb S^2)=2<\mu_4(\mathbb S^2)=6.\] But \(\Omega^\varepsilon=\mathbb S^2\setminus\psi(\varepsilon\omega)\) has two connected components: the image of the small interior disk and the exterior of the outer circle. Constants on these components form a two-dimensional Neumann kernel. With the paper’s multiplicity indexing and its mean-zero variational definition on page 1, \[
 \mu_0(\Omega^\varepsilon)=\mu_1(\Omega^\varepsilon)=0.\] This contradicts the claimed convergence of the index \(i=1\) to \(2\).\par
Equation reference: kappa. The proposed coefficient can also be calculated without any ambiguity from the constants in the vector corrector. On the complementary components take \[W(x)=9\frac{x-c}{|x-c|^2}\quad(r>3),\qquad
 W(x)=-(x-c)\quad(r<1).\] Both expressions are harmonic and satisfy \(\partial_{\nu_\omega}W=-\nu_\omega\), and the outer expression decays. An arbitrary constant on the inner component shows that the uniqueness claim on page 4 fails, but all such choices have the same gradients. For each coordinate, the outer and inner Dirichlet energies are \(9\pi\) and \(\pi\); the mixed energy is zero. As \(|\omega|=8\pi\), one obtains \(M_\omega=18\pi I_2\). For the normalized coordinate eigenfunctions \(u_i=\sqrt{3/(4\pi)}\,x_i\) on the sphere, the perturbation matrix is \[
 \frac{3}{4\pi}
 \begin{pmatrix}
 -18\pi&0&0\\0&-18\pi&0\\0&0&16\pi
 \end{pmatrix}
 =\operatorname{diag}\left(-\frac{27}{2},-\frac{27}{2},12\right).\] Proposition 3 would require \[\left|0-2+\frac{27}{2}\varepsilon^2\right|\leq C\varepsilon^{9/4},\] which is impossible as \(\varepsilon\downarrow0\). This is a fixed admissible geometry and an entire limiting family, not a single numerical comparison.\par
The matrix corrector \(L\) on page 4 faces a further direct obstruction. On \(D_1(c)\) its boundary condition would imply, for either diagonal entry and the outward disk normal \(n\), \[\int_{\partial D_1(c)}\partial_n L_{ii}\,\mathrm ds
 =\int_{\partial D_1(c)}x_i n_i\,\mathrm ds=|D_1(c)|=\pi.\] A harmonic function on the disk has zero total outward flux. Thus this component of the proposed exterior problem is incompatible. Summing a flux identity over all boundary components does not remove the compatibility requirement on each bounded component.\par
\subsection*{Counterreadings, impact, and repair limits}
If hole is intended to mean a filled obstacle with connected exterior, that interpretation adds the missing hypothesis. It does not follow from smoothly bounded closed set, and the annulus itself is connected. The additional zero mode cannot be discarded silently: the manuscript counts multiplicities and fixes indices. The compactness argument on page 7 fails for the same reason: the \(L^2\)-normalized indicator of the shrinking interior component has zero Dirichlet energy and no strongly convergent \(L^2\) subsequence after extension by zero.\par
A suitable local correction is to require each removed set to have connected complement, for example to be the closure of a smooth bounded simply connected domain, and to add the matching assumption in Lemma 7. After inversion the complement is then a connected bounded domain. Subtracting one mean permits Poincaré’s inequality; a trace/extension estimate and harmonic energy minimization give the needed Dirichlet-energy bound. Two-dimensional conformal invariance supports the scaling. This identifies the failed step and a correction compatible with the actual applications. It is not an amended theorem allowing cavities, which would need componentwise constants and different spectral indexing.\par
Section 4 uses filled ellipses; Section 6 explicitly assumes a fixed simply connected hole. The counterexample lies outside both families after the stated correction. No demonstrated dependence failure survives in those applications. This local scope failure is therefore classified as Minor, without claiming a refutation of Theorem 1.\par
\section*{Further checks and evidence limits}
\subsection*{Analytical versus numerical assertions}
Section 5 explicitly identifies its helmet-domain arguments as non-rigorous. The curves in Figure 3 and the mode transition in Figure 4 were actually observed; their numerical intervals are not promoted here to certified bounds. The graph-limit expression and its sign mechanism were examined as heuristics. A rigorous theorem for the lower numerical threshold is neither claimed there nor established by this review.\par
In Section 6 the sphere eigenspace matrix gives \[\mu_1(\Omega^\varepsilon)+\mu_2(\Omega^\varepsilon)
 =4-3\,\operatorname{cap}[\omega]^2\varepsilon^2+o(\varepsilon^2).\] The argument uses fixed smooth simply connected shapes and the area–capacity comparison with the disk. It expressly does not optimize over arbitrary shape-changing sequences of large domains. The difference between fixed planar hole area and the corresponding spherical area changes only higher orders in the comparison. No counterexample was obtained within this scope.\par
\subsection*{Independent calculation and static code inspection}
The decisive arguments [energy] (\#energy)–[kappa] (\#kappa) are analytical. A reviewer-authored script checked the rational coefficients, the origin-in-annulus condition, and an independent composite-Simpson evaluation of [energy] (\#energy). It returned \(5.7192017347602535\) for the nonzero energy and zero exterior energy. This is an arithmetic and quadrature check, not an eigenvalue solver or a replacement for the topological and PDE proofs.\par
The preserved author repository is pinned to commit `159ee85f54e966c530ae5de0ceb093b5c6f7d56e`. Its README and two FreeFem++ scripts were read statically. The eigenvalue script assembles surface \(P_1\) stiffness and mass matrices; the mesh-generation script describes punctured caps. Four retained value tables were programmatically indexed, not reproduced as scientific computations. The 471 mesh files in the preserved index were not acquired, and a direct mapping from the retained material to the helmet curves in Figure 3 was not established. No author code was executed. These are reproducibility limits, not confirmed mathematical defects or evidence that the published curves are false.\par
\section*{Assessment}
Item | Status | Demonstrated consequence\par
Problem 1 | Minor | Lemma 7 and Proposition 3 need a connected-exterior hypothesis; a cavity also obstructs \(L\).\par
Theorem 1 | No further confirmed defect in the examined deduction | The filled-ellipse argument avoids the example; Proposition 2 remains an explicitly identified external input.\par
Section 5 | Numerical; heuristic | Figures were inspected; numerical reproduction and certified thresholds are not claimed.\par
An explicit topological restriction and a matching correction of the extension/exterior-problem discussion are requested. The absence of additional confirmed findings is limited to the work and evidence described above.\par
\subsection*{Supporting checks}
No substantive prior report for this paper was recovered. Preserved source bytes and the pinned repository snapshot were reused. The reading, mathematical arguments, reviewer calculation, report, compilation and actual image inspection belong to the current scientific conversation.\par
The scientific runtime self-description is GPT-6 Astra Pro; backend equivalence was not independently established. The controller contributed scheduling and byte collection.\par
The report uses the existing Tectonic 0.17.0 runtime without shell escape. Exact source, TeX, compiler, per-paper dependency, log, PDF and actually viewed image bindings are recorded alongside it. Delivery completion records initial-review work only. The scientific decision remains pending, human reviews are empty, and independent scientific acceptance is reserved for a later reviewer.\par
9 D. Bucur, R. S. Laugesen, E. Martinet, and M. Nahon, Spherical Caps do Not Always Maximize Neumann Eigenvalues on the Sphere, arXiv:2503.15385v1, 19 March 2025. Frozen journal DOI: `10.1007/s00039-025-00721-1`. J. J. Langford and R. S. Laugesen, Maximizers beyond the hemisphere for the second Neumann eigenvalue, Math. Ann. 386 (2023), 2255–2281. DOI: `10.1007/s00208-022-02455-z`. Only accessible publisher material and appendices were consulted; the principal propositions were not fully independently audited. E. Martinet, PuncturedSpheres, preserved repository snapshot, commit `159ee85f54e966c530ae5de0ceb093b5c6f7d56e`. Static inspection only; no author code execution.\par
\end{document}
